\chapter{Profesor invitado procesos de selección}

En esta sesión contamos con la participación de un profesional con una larga trayectoria en el ámbito de los recursos humanos, responsable de la dirección de talento en una multinacional del sector industrial. Su intervención ofreció una mirada muy completa sobre lo que ocurre “al otro lado de la mesa” en un proceso de selección, permitiéndonos conectar los contenidos teóricos vistos en clase con la realidad práctica de las empresas. 

La sesión comenzó con una explicación sobre la magnitud y complejidad de los procesos de selección en una gran organización, que puede gestionar cientos de incorporaciones al año en distintos países y perfiles. Esta introducción permitió comprender la importancia de la planificación de personal o \textit{workforce planning}, donde se definen las necesidades futuras de plantilla en función de los proyectos y líneas de negocio. El ponente insistió en que, sin personas, no hay proyecto, subrayando el papel estratégico de la función de selección. Además, destacó el contexto actual de escasez de talento, en el que muchas compañías compiten por perfiles similares, lo que convierte la atracción y retención de profesionales en un reto compartido por todo el sector.  

Uno de los puntos más interesantes fue el contraste entre las percepciones del candidato y del seleccionador. Como estudiantes y futuros profesionales, solemos asociar la entrevista con nervios o inseguridad, sin reparar en la presión y la responsabilidad que también siente la persona que selecciona. En este sentido, se remarcó que la entrevista no debe verse como una prueba unidireccional, sino como una conversación entre dos partes que buscan un encaje mutuo. Esta visión coincide con lo trabajado en clase sobre el contrato psicológico y la necesidad de generar relaciones laborales basadas en la confianza y la transparencia.  

A partir de ahí, el ponente nos guió a través de las principales fases del proceso de selección. En la primera, la inscripción, insistió en la importancia de tener el perfil de \textit{LinkedIn} y el currículum siempre actualizados, ya que las oportunidades pueden surgir en cualquier momento. En la segunda fase, centrada en el diseño del CV, destacó que este es nuestro escaparate: solo disponemos de unos segundos para captar la atención de quien lo revisa. Por ello, debe ser claro, ordenado, visualmente limpio y sin errores. Defendió también que incluir una fotografía puede ser positivo si transmite profesionalidad y refuerza la candidatura, siempre que se utilice con criterio.  

Durante la fase de entrevista, se realizó un \textit{role play} que ilustró con claridad la diferencia entre un candidato preparado y uno que no lo está. La primera candidata, que había investigado sobre la empresa y mostraba interés genuino, transmitía profesionalidad y autoconfianza; mientras que el segundo, desinformado y con un tono poco adecuado, proyectaba desinterés. La actividad sirvió para reflexionar sobre la importancia de la preparación y la actitud. En clase se había mencionado que el éxito de una entrevista depende en gran parte de la autoconciencia y la capacidad de comunicación, y esta práctica lo demostró de forma muy visual.  

Entre los consejos más valiosos, el ponente insistió en que la entrevista comienza mucho antes de sentarse frente al entrevistador. La preparación previa supone, según él, hasta un 70\% del éxito. Es necesario investigar sobre la empresa, definir los mensajes clave que se quieren transmitir y reflexionar sobre qué impresión se desea causar. También mencionó la importancia de la puntualidad, la coherencia entre lenguaje verbal y no verbal, y la actitud positiva. Recomendó mostrar interés y entusiasmo, pero con naturalidad, evitando sobreactuar. La entrevista, dijo, “es un diálogo entre iguales donde ambas partes tienen algo que ganar”.  

Asimismo, la exposición incluyó una parte muy interesante sobre la digitalización de los procesos de selección. Se presentaron herramientas como los sistemas de seguimiento de candidatos (\textit{Applicant Tracking Systems}) y plataformas como \textit{Workday}, que automatizan parte de la comunicación y mejoran la experiencia del candidato. También se habló del papel incipiente de la inteligencia artificial en la fase de cribado, especialmente en procesos masivos. Sin embargo, se advirtió que en perfiles técnicos y de ingeniería sigue siendo esencial la valoración humana, tanto por la complejidad del puesto como por la necesidad de evaluar la adecuación cultural.  

Otro aspecto destacado fue el concepto de \textbf{experiencia del candidato}. Del mismo modo que los candidatos se esfuerzan por resultar atractivos, las empresas también deben cuidar su imagen y ofrecer procesos claros, ágiles y comunicativos. En un mercado competitivo, las compañías no solo buscan talento, sino que también deben “venderse” como buenos lugares para trabajar. Esta idea enlaza directamente con lo visto en clase sobre marca empleadora y gestión de la reputación corporativa.  

La sesión concluyó con una invitación a la reflexión personal. Prepararse para un proceso de selección no se limita a aprender respuestas estándar, sino a conocerse a uno mismo, identificar los propios valores y definir qué tipo de entorno profesional se busca. También se insistió en la importancia de la sinceridad: tanto empresas como candidatos deben ser honestos durante la entrevista, porque un encaje mal planteado termina siendo negativo para ambas partes.  

En definitiva, la sesión aportó una visión muy práctica sobre los procesos de selección y permitió comprender cómo las decisiones en materia de talento combinan análisis técnico, intuición y humanidad. Personalmente, me resultó especialmente útil entender que la entrevista no es un examen, sino una oportunidad para generar confianza mutua y mostrar coherencia entre lo que uno dice y lo que transmite. Me quedo con la idea de que la preparación, la autenticidad y el interés genuino son las claves para destacar en un entorno profesional cada vez más exigente y cambiante. 
